\documentclass[11pt]{article}
\usepackage[a4paper, margin=1in]{geometry}
\usepackage{amsmath}
\usepackage{amssymb}
\usepackage{enumitem}
\usepackage{fancyvrb}
\usepackage{xcolor}
\usepackage{hyperref}
\usepackage{booktabs}
\usepackage{array}
\usepackage{longtable}
\usepackage{graphicx}
\usepackage{fancyhdr}

\setlength{\parindent}{0pt}
\setlength{\parskip}{0.8em}

\pagestyle{fancy}
\fancyhf{}
\cfoot{\thepage}
\renewcommand{\headrulewidth}{0pt}
\renewcommand{\footrulewidth}{0pt}

\newcommand{\code}[1]{\texttt{\small #1}}

\begin{document}

\begin{center}
    \vspace*{5cm}
    {\LARGE \textbf{Michael Polley}} \\[2em]
    \normalsize Student ID: T00774397 \\[2em]
    NFL Big Data Bowl 2027 \\[2em]
    Research Report: Beyond the Stopwatch, Behind the Draft Board
    \vfill
\end{center}
\thispagestyle{empty}
\newpage

\noindent\textbf{1.}
\noindent\textbf{(a)}
We consider whether 10 Hz optical Combine tracking predicts NFL careers beyond what teams already know on draft day. The cohort is 510 prospects drafted or signed between 2023 and 2025. A stopwatch records only the elapsed time
\[ T = \int_{0}^{40} \frac{dx}{v(x)}, \]
while tracking records the whole speed curve $s(t) = \|\mathbf{v}(t)\|$ and the acceleration magnitude $a(t)$ in every drill. NFL playing time, however, is assigned largely by draft slot. Thus the question a scout faces is not whether a metric correlates with snaps but whether it predicts snaps \emph{conditional on} draft position.

\noindent\textbf{(b)}
We fit every model twice. The raw specification regresses an outcome on the Combine metric and the original covariate (40-yard dash time, or weight for defensive linemen). The controlled specification adds $\log(\text{pick})$, an undrafted indicator (undrafted players coded as pick 262), draft-class fixed effects for seasons of exposure, and position fixed effects for DL, DB and OL:
\[ y_i = \beta_0 + \beta_1 m_i + \beta_2 c_i + \delta \log(\text{pick}_i) + \eta\,\text{UDFA}_i + \lambda_{\text{class}(i)} + \epsilon_i . \]
Ordinary least squares models use heteroskedasticity-robust (HC1) standard errors. Starter status ($\ge 10$ career starts) uses a logit. Efficiency outcomes use regular-season snaps only (249,269 of 314,197 rows; the remainder are preseason and playoff snaps).

\noindent\textbf{2.}
\noindent\textbf{(a)}
We define the metrics. For receivers $s_{\text{cut}}$ is the median speed over every frame of the curl, comeback, speed-out and slant drills. Prospects did not run the same drills, so we also define the drill-standardised score
\[ z_{\text{cut},i} = \frac{1}{|D_i|} \sum_{d \in D_i} \frac{\tilde s_{i,d} - \bar s_d}{\sigma_d}, \]
where $\tilde s_{i,d}$ is player $i$'s median speed in drill $d$ and $D_i$ is the set of cutting drills he ran. For linemen, $\gamma$ is the 90th percentile of $a$, and Gamma-Over-Expected is the residual on weight:
\[ \gamma_{\text{OE},i} = \gamma_i - \big(\hat\alpha + \hat\beta\,\text{Weight}_i\big). \]
Weight explains 29\% of the variance in $\gamma$ ($\hat\beta = -0.019$ km/h/s per lb). The column $a$ is a magnitude, so $\gamma$ counts hard braking as well as burst. Pooled over all drills, $\gamma$ also correlates $0.36$ with the number of drills run. Thus we also compute $\gamma$ within the pass-rush drill, which all 118 linemen ran.

\noindent\textbf{(b)}
We describe regular-season game speed across 18.1 million tracking frames. The median speed is 2.06 yd/s (6.8 km/h) and the 95th percentile is 6.66 yd/s (21.9 km/h). Of all frames 66.3\% fall below 10 km/h, 32.2\% fall between 10 and 26 km/h, and 1.5\% exceed 26 km/h. At the snap level, 7.0\% of player-snaps reach 26 km/h at some point. Thus most football is played at moderate speed, and movement through a cut is a reasonable target for measurement.

\newpage

\noindent\textbf{3.}
\noindent\textbf{(a)}
We test the four metrics behind the original claims, one per position group.

\begin{center}
\small
\begin{tabular}{lcc}
\toprule
\textbf{Metric $\to$ outcome} & \textbf{Metric alone} & \textbf{+ draft controls} \\
\midrule
WR $s_{\text{cut}}$ $\to$ career snaps (per km/h), $N=96$ & $+114$, $p = 0.002$ & $-10$, $p = 0.74$ \\
WR $s_{\text{cut}}$ $\to$ starter (odds ratio), $N=96$ & $1.63$, $p = 0.026$ & $1.11$, $p = 0.76$ \\
DL $\gamma_{\text{OE}}$ (all drills) $\to$ starter (odds ratio), $N=118$ & $1.85$, $p = 0.042$ & $1.52$, $p = 0.35$ \\
DB break speed $\to$ career snaps (per km/h), $N=101$ & $+192$, $p = 0.011$ & $-15$, $p = 0.80$ \\
OL $\gamma_{\text{OE}}$ $\to$ career snaps, $N=95$ & $+155$, $p = 0.015$ & $+58$, $p = 0.32$ \\
\bottomrule
\end{tabular}
\end{center}

Every raw association is significant at 5\%, and every controlled association is not. For receiver snaps the draft variables alone give $R^2 = 0.511$; adding $s_{\text{cut}}$ and the 40 time gives $R^2 = 0.513$. Receivers with higher $s_{\text{cut}}$ are drafted earlier (slope on $\log(\text{pick})$, $p = 0.012$). Thus the Combine metric and the draft slot carry the same information, and the metric adds none.

\noindent\textbf{(b)}
We partition the 106 receivers by career starts and compare pooled cut speed with draft capital.

\begin{center}
\small
\begin{tabular}{lcccc}
\toprule
\textbf{Tier} & \textbf{N} & \textbf{Mean $s_{\text{cut}}$} & \textbf{Mean pick (UDFA $=262$)} & \textbf{Undrafted} \\
\midrule
Top ($\ge 20$ starts) & 14 & $18.03$ km/h & 62 & 0\% \\
Middle (3--19) & 40 & $17.19$ km/h & 116 & 10\% \\
Fringe ($\le 2$) & 52 & $16.53$ km/h & 184 & 50\% \\
\bottomrule
\end{tabular}
\end{center}

Welch's two-sample test comparing Top and Fringe gives $t = 4.75$, $p < 0.001$, with Cohen's $d = 1.00$. Half of the Fringe tier went undrafted, and the Top tier averaged pick 62. Thus the raw gap in cut speed reflects who was drafted, not who outplayed their slot.

\noindent\textbf{(c)}
We account for multiple testing. Across all 22 controlled tests (six metrics; volume and efficiency outcomes), one has $p < 0.05$: DL $\gamma_{\text{OE}}$ (all drills) on career snaps, $+85$ snaps per km/h/s, $p = 0.048$. About one such result is expected by chance. Thus after Holm's correction no controlled effect is significant.

\newpage

\noindent\textbf{4.}
\noindent\textbf{(a)}
We test per-snap efficiency, which depends less on draft investment than snap counts do. Receivers need at least 100 regular-season routes ($N = 57$ with a 40 time) and linemen at least 100 pass rushes ($N = 69$). With draft controls, the drill-standardised score $z_{\text{cut}}$ raises yards per route run by $0.32$ per unit ($p = 0.058$); pooled $s_{\text{cut}}$ shows no effect ($p = 0.56$). Pass-rush-drill $\gamma_{\text{OE}}$ raises starter odds by a factor of $1.60$ ($p = 0.059$) but shows no link to pressure rate ($p = 0.55$) or get-off time ($p = 0.32$). Separation at the throw is not higher for faster cutters: the raw slope on $s_{\text{cut}}$ is $-0.08$ yd per km/h ($p = 0.013$), and $p = 0.16$ with controls. Thus two leads remain near the 5\% line, and the cutting-separation mechanism is not supported.

\noindent\textbf{(b)}
We rank every available predictor: draft slot, Combine results, NGS scores and the tracking metrics. For starter status we use cross-validated AUC (five folds, twenty repeats). Draft pick ranks first at every position (AUC 0.87 WR, 0.82 DL, 0.90 DB, 0.93 OL), followed by the NGS final and college-production scores. No predictor raises AUC by more than 0.01 over draft pick, apart from two small-sample DB cases. For efficiency we use rank correlations after removing draft pick (and edge versus interior for linemen). For receiver yards per route run the strongest is $z_{\text{cut}}$ ($+0.33$, $p = 0.008$). For pass-rush pressure rate the strongest are the NGS final score ($+0.36$, $p = 0.003$), the 10-yard split ($-0.29$, $p = 0.034$) and the 40 time ($-0.28$, $p = 0.038$). Seven of these 30 correlations fall below $p = 0.05$ where 1.5 are expected, but none survives Holm's correction. Thus Combine information appears in per-snap efficiency, not in playing time, and for pass rushers it includes the stopwatch.

\noindent\textbf{(c)}
We examine one comeback-route rep each for Puka Nacua (pick 177, 201 lb) and Quentin Johnston (pick 21, 208 lb). Both reach a stem peak near 27 km/h (27.6 and 27.1) and slow to about 12 km/h (12.2 and 12.4). Each sheds about 385--390 kg$\cdot$m/s of momentum, with peak braking of $-5.77$ and $-5.23$ yd/s$^2$. One second after his slowest frame Nacua is back at 22.7 km/h, while Johnston is at 15.7 km/h and slowing. Nacua's path breaks back toward the line; Johnston's turns closer to $90^\circ$. Thus re-acceleration out of the break is a candidate feature, but one rep per player is an illustration, not evidence.

\noindent\textbf{5.}
\noindent\textbf{(a)}
We summarise the guidance for front offices.

\begin{center}
\small
\setlength{\tabcolsep}{4pt}
\begin{tabular}{p{0.32\textwidth}p{0.6\textwidth}}
\toprule
\textbf{Practice} & \textbf{Reason} \\
\midrule
Test metrics conditional on draft slot & Draft variables explain 51\% of receiver snap variance; tracking adds 0.2 points \\
Standardise within each drill & Pooled medians mix drills; Smith-Njigba's pooled median is 24.0 km/h, his cutting median 19.6 \\
Judge on per-snap efficiency & Snaps and starts follow draft investment \\
Keep the stopwatch for pass rushers & 10-yard split relates to pressure rate at a given pick and position ($p = 0.034$) \\
Pre-register the leads for 2026 & $z_{\text{cut}} \to$ YPRR; 10-yard split $\to$ pressure rate; pass-rush $\gamma_{\text{OE}} \to$ starter \\
\bottomrule
\end{tabular}
\end{center}

\noindent\textbf{(b)}
We state the limitations. Careers span one to three seasons, efficiency samples are small, each drill has one tracked rep per player, and $a$ cannot separate acceleration from braking. Thus Combine tracking has not yet been shown to add information beyond the draft board. A pre-registered test on the 2026 class is the next step, with the leads defined in advance.

\end{document}
